\section{Parasitics}
\label{sec:theory_general_parasitics}

In real-world circuits, components are not ideal and exhibit
parasitic characteristics that can affect circuit performance.
This is because every material has some inherent inductance
and resistance while nearby conductors can introduce parasitic
capacitance.

\subsection{Types of Parasitics}
\begin{itemize}
  \item \textbf{Parasitic Capacitance:} Unintended capacitance
    between conductive parts of a circuit, such as between
    traces on a PCB or between component leads. This can lead to
    signal coupling and affect high-frequency performance.
  \item \textbf{Parasitic Inductance:} Unintended inductance
    in circuit elements, often due to the physical layout of
    components and wiring. This can cause voltage spikes and
    affect switching performance in high-speed circuits.
  \item \textbf{Parasitic Resistance:} Unintended resistance
    in components and connections, which can lead to power
    losses and voltage drops.
\end{itemize}

\subsection{Resistor Parasitics}
Below I've shown the model of a real resistor including its
parasitic elements.
\begin{center}
  \begin{circuitikz}
    % Paths, nodes and wires:
	  \draw (8.75, 7) to[european resistor, l={$R_2$}] (10.5, 7);
	  \draw (8.75, 7) to[european inductor, l_={$L_2$}] (7, 7);
	  \draw (4.5, 7) to[european resistor, l={$R_1$}] (6.25, 7);
	  \draw (4.5, 7) to[european inductor, l_={$L_1$}] (2.75, 7);
	  \draw (13, 7) to[european resistor, l={$R_3$}] (14.75, 7);
	  \draw (13, 7) to[european inductor, l_={$L_3$}] (11.25, 7);
	  \draw (8, 8.75) to[capacitor, l={$C_{13}$}] (9.5, 8.75);
	  \draw (8, 8.75) -- (2.75, 8.75) -- (2.75, 7);
	  \draw (6.25, 7) -- (7, 7);
	  \draw (10.5, 7) -- (11.25, 7);
	  \draw (14.75, 7) |- (9.5, 8.75);
	  \node[shape=rectangle, draw, line width=3pt, dash pattern={on 3pt off 6pt}, minimum width=4.394cm, minimum height=1.644cm] at (8.75, 7.125){};
	  \node[shape=rectangle, draw, line width=3pt, dash pattern={on 3pt off 6pt}, minimum width=3.644cm, minimum height=1.144cm] at (4.625, 7.125){};
	  \node[shape=rectangle, draw, line width=3pt, dash pattern={on 3pt off 6pt}, minimum width=3.644cm, minimum height=1.144cm] at (12.75, 7.125){};
  \end{circuitikz}
\end{center}
